% Source: https://artofproblemsolving.com/community/c3h2795007_2011_tx_mmc__texas__mathworks_math_contest
% Stable identity: aops-topic-2795007-q04
Put the corner at the origin, so the circle has centre $(r,r)$ and radius $r$. The far vertex of the rectangle is $V=(6,3)$.
\begin{asy}
usepackage("amsmath");
size(5cm);
defaultpen(fontsize(10pt));
real r=15;
draw((0,0)--(30,0)--(30,30)--(0,30)--cycle);draw(circle((15,15),15));
filldraw((0,0)--(6,0)--(6,3)--(0,3)--cycle,gray(0.8));
label("$6$",(3,0),S);label("$3$",(0,1.5),W);

pair O=(15,15),V=(6,3);draw(O--V,dashed);dot(O);dot(V);label("$O$",O,NE);label("$V$",V,NW);draw(O--(15,3),dotted);draw((6,3)--(15,3),dotted);label("$r-6$",(12,3),N);label("$r-3$",(15,9),E);
\end{asy}

Then
\[(r-6)^2+(r-3)^2=r^2\iff r^2-18r+45=0\iff (r-3)(r-15)=0.\]
If $r=3$ the square has side $6$ and the rectangle covers half of it, containing points inside the circle, which is not allowed. Hence $r=\boxed{15}$ cm.
